\section*{第 \ref{ch:g3:sound-around-us} 章——身边的声音}

\begin{solution}{exo:g3:sound-around-us:1}
非常快速的来回颤动。吉他：被拨动的弦。
鼓：被敲击的鼓面。人声：喉咙里的两片小声带。
苍蝇：拍动的翅膀。
\end{solution}

\begin{solution}{exo:g3:sound-around-us:2}
发声的鼓面就是在颤动的鼓面：米粒跳起，是因为
它下面的鼓面在快速地上下运动——哪怕只是
一声喊叫而不是触碰，就让它\omterm{def:g3:sound-around-us:vibration}{振动}起来。
\end{solution}

\begin{solution}{exo:g3:sound-around-us:3}
手握住铙钹，就让它停止颤动，没有\omterm{def:g3:sound-around-us:vibration}{振动}，
就没有\omterm{def:g3:sound-around-us:sound}{声音}——这就是规律“无\omterm{def:g3:sound-around-us:vibration}{振动}，无\omterm{def:g3:sound-around-us:sound}{声音}”。
\end{solution}

\begin{solution}{exo:g3:sound-around-us:4}
拨得更用力：皮筋摆动得更宽——模糊的范围更大——
嘣声也更响。颤动越大，\omterm{def:g3:sound-around-us:sound}{声音}越大。
\end{solution}

\begin{solution}{exo:g3:sound-around-us:5}
短而被捏住的皮筋；小笛子；最细的弦。
短的、小的、细的东西颤动得更快，\omterm{def:g3:sound-around-us:sound}{声音}更高。
\end{solution}

\begin{solution}{exo:g3:sound-around-us:6}
他们感到嗡嗡作响——那是喉咙里的声带在\omterm{def:g3:sound-around-us:vibration}{振动}。嗡嗡声
一旦哼唱停止就立即消失：没有\omterm{def:g3:sound-around-us:vibration}{振动}，就没有\omterm{def:g3:sound-around-us:sound}{声音}。
\end{solution}

\begin{solution}{exo:g3:sound-around-us:7}
拍手使\omterm{def:g2:air-around-us:air}{空气}颤动；颤动像池塘中的涟漪一样
在\omterm{def:g2:air-around-us:air}{空气}中扩散，传到耳朵，使耳朵的小鼓膜
颤动起来——于是听到了拍手声。
\end{solution}

\begin{solution}{exo:g3:sound-around-us:8}
很响的\omterm{def:g3:sound-around-us:sound}{声音}会震动耳朵里的小鼓膜——那是耳朵自己的
米粒大小的“鼓”，比纸还薄——强烈的震动会永久损伤它，
且无法再长回来。
\end{solution}

\begin{solution}{exo:g3:sound-around-us:9}
短的齿片发出高音，长的齿片发出低音：
短的东西颤动得更快，\omterm{def:g3:sound-around-us:sound}{声音}更高。
\end{solution}

\begin{solution}{exo:g3:sound-around-us:10}
瓶内的\omterm{def:g2:air-around-us:air}{空气}是\omterm{def:g3:sound-around-us:vibration}{振动}体——你的呼气使它颤抖。
加水后，水面上方的\omterm{def:g2:air-around-us:air}{空气}柱变得\emph{较短}，
短的\omterm{def:g3:sound-around-us:sound}{声音}就高，所以哨声会升高。如果
玻璃瓶本身是发声体，加水几乎不会改变
音调。
\end{solution}

\begin{solution}{exo:g3:sound-around-us:11}
\omterm{def:g3:sound-around-us:sound}{声音}需要\emph{\omterm{def:g3:sound-around-us:vibration}{振动}}：一种来回的颤动，
必须足够快，才能使\omterm{def:g2:air-around-us:air}{空气}\omterm{def:g3:sound-around-us:vibration}{振动}，就像让米粒跳舞的那张鼓面。
无声的屏幕上画面在变化，但没有任何东西使\omterm{def:g2:air-around-us:air}{空气}颤动——
颜色的变化不会震动任何东西，
因此耳朵无从听到。
\end{solution}
